\chapter{Introdução}
\label{cap:introducao}

\section{Motivação}
\label{sec:motivacao}

A computação em nuvem consolidou-se como a principal forma de disponibilizar infraestrutura computacional para empresas e instituições. Nesse modelo, o provedor mantém grandes centros de dados e precisa decidir, de forma contínua, quanto de processamento, memória e armazenamento deve estar disponível para atender à demanda de seus clientes. Como a demanda tende a crescer, impulsionada sobretudo por aplicações de Inteligência Artificial e de análise de grandes volumes de dados, a gestão eficiente desses recursos tornou-se um fator central para a sustentabilidade econômica dos provedores \cite{rossi2023forecasting}.

A dificuldade está no fato de que a carga de trabalho não é constante. Ela apresenta picos, padrões periódicos e variações aparentemente aleatórias. Diante disso, a abordagem mais comum é a alocação reativa, na qual o sistema ajusta os recursos apenas depois de observar uma mudança na demanda, normalmente por meio de regras de limiar sobre a utilização. Como a reação ocorre com atraso, esse tipo de estratégia tende a degradar o desempenho durante os picos e a desperdiçar recursos nos períodos de baixa utilização \cite{bodra2025ml}.

Esse cenário cria um dilema permanente. Manter recursos em excesso (superprovisionamento) garante a qualidade do serviço, mas eleva o custo operacional. Manter recursos de menos (subprovisionamento) reduz o custo, mas provoca violações dos Acordos de Nível de Serviço (\textit{Service Level Agreements}, SLA) firmados com os clientes. Uma alternativa é a alocação preditiva, que utiliza o histórico de uso para antecipar a demanda futura e provisionar recursos antes que eles sejam necessários \cite{bodra2025ml}. Nesse contexto, prever a carga com precisão passa a ser uma pré-condição para otimizar a alocação \cite{rossi2023forecasting}.

A previsão de carga em nuvem foi tradicionalmente tratada com modelos estatísticos lineares, como o modelo autorregressivo integrado de médias móveis (ARIMA) \cite{calheiros2015arima}. Mais recentemente, redes neurais profundas, em especial as arquiteturas recorrentes, passaram a ser aplicadas ao problema com bons resultados, inclusive em abordagens multivariadas que consideram simultaneamente o uso de processador e de memória \cite{xu2022esdnn}.

\section{Definição do problema}
\label{sec:definicao-problema}

O problema tratado neste trabalho tem duas partes encadeadas. A primeira é um problema de previsão: dado o histórico de utilização de recursos de um \textit{cluster}, amostrado em intervalos regulares, estimar a demanda de processador e de memória em um horizonte curto. A segunda é um problema de decisão: a partir dessa estimativa, determinar quantas instâncias de recursos (máquinas virtuais ou contêineres) devem permanecer ativas em cada intervalo.

De forma preliminar, a decisão pode ser formulada como a minimização do custo operacional sujeita a uma restrição de qualidade de serviço:

\begin{equation}
\label{eq:problema}
\min_{n_1, \dots, n_T} \; \sum_{t=1}^{T} c \cdot n_t
\quad \text{sujeito a} \quad
\frac{1}{T} \sum_{t=1}^{T} \mathbf{1}\left[\, d_t > k \cdot n_t \,\right] \leq \varepsilon,
\end{equation}

\noindent em que $n_t$ é o número de instâncias ativas no intervalo $t$, $c$ é o custo de uma instância por intervalo, $k$ é a capacidade de uma instância, $d_t$ é a demanda real observada, $\mathbf{1}[\cdot]$ é a função indicadora e $\varepsilon$ é a fração máxima de intervalos com violação de SLA tolerada. Como $d_t$ não é conhecida no momento da decisão, a escolha de $n_t$ depende da previsão $\hat{d}_t$, e a qualidade dessa previsão afeta diretamente o custo e o número de violações. Essa formulação é refinada no Capítulo~\ref{cap:objetivos}, que define como a restrição é verificada na prática (Equação~\ref{eq:calibracao}).

\section{Justificativa}
\label{sec:justificativa}

A relevância do problema é econômica e operacional. Cada ponto percentual de recursos ociosos representa custo para o provedor, e cada violação de SLA representa perda de qualidade percebida pelo cliente e, eventualmente, penalidades contratuais. Apesar do volume de trabalhos sobre previsão de carga, falta uma avaliação de ponta a ponta, com dados públicos, que meça quanto a previsão reduz o custo de alocação quando redes neurais e métodos tradicionais são submetidos à mesma política, lacuna discutida na Seção~\ref{sec:lacuna}.

\section{Escopo e delimitação}
\label{sec:escopo}

\textbf{Incluído.} Rastros públicos da Alibaba \cite{alibaba2018trace} e do Google \cite{wilkes2020google}, agregados em séries de uso de processador e de memória; horizontes de previsão curtos, na escala de minutos; alocação em um centro de dados simulado, com custo por instância ativa e violação de SLA nos intervalos em que a demanda excede a capacidade; e comparação entre uma rede neural e métodos tradicionais de previsão sob a mesma política de alocação, ao lado das políticas de referência da Seção~\ref{sec:solucao}.

\textbf{Fora do escopo.} Implantação em uma nuvem real, posicionamento e migração de máquinas virtuais entre servidores físicos, consumo de energia, múltiplos provedores e aprendizado por reforço para a decisão. Na simulação, o posicionamento segue a política padrão do simulador, a mesma para todos os previsores.

\section{Organização do documento}
\label{sec:organizacao}

O restante deste documento está organizado da seguinte forma. O Capítulo~\ref{cap:situacao-atual} apresenta a terminologia básica e as principais abordagens existentes para previsão de carga e alocação de recursos em nuvem, com suas limitações, além das tecnologias utilizadas. O Capítulo~\ref{cap:objetivos} descreve a solução proposta, seus objetivos específicos e o que será efetivamente implementado. O Capítulo~\ref{cap:plano-acao} revisa o plano de ação, comparando as atividades previstas com as realizadas, e apresenta os resultados preliminares do protótipo. Por fim, o Capítulo~\ref{cap:cronogramas} apresenta os cronogramas do Projeto Final I (PF~I) e do Projeto Final II (PF~II).